\documentclass[11pt,a4paper]{article}
\usepackage{amsmath,amssymb,amsthm}
\usepackage[margin=2.5cm]{geometry}
\usepackage{hyperref}

\newtheorem{definition}{Definition}[section]
\newtheorem{theorem}[definition]{Theorem}
\newtheorem{proposition}[definition]{Proposition}
\newtheorem{lemma}[definition]{Lemma}
\newtheorem{corollary}[definition]{Corollary}
\newtheorem{conjecture}[definition]{Conjecture}
\newtheorem{remark}[definition]{Remark}

\title{Hyperseed Formalization:\\
  What Should an AGI Curriculum Look Like?}
\author{ProtomegaTron (automated formalization)}
\date{2026-09-17}

\begin{document}
\maketitle

\begin{abstract}
We formalize the intellectual structure of Ben Goertzel's Substack article
``What Should an AGI Curriculum Look Like, at the Dawn of the AGI Era?''
(2026-07-23) within the Hyperseed ontology. The article announces the first
accredited AGI degree program (MS in AGI at CIHS) and articulates what AGI
education should contain. Each structural insight maps onto Hyperseed
structures: the three curriculum levels (Mind, Brain, Experience) as fiber
decomposition, vibe coding vs.\ deep thinking as 0-cell surface optimization
vs.\ directed higher-cell reasoning, ethics as value possession rather than
rule adoption, attention as parallel transport budget, and multiple
perspectives as a sheaf of local sections.
\end{abstract}

\tableofcontents

%% =================================================================
\section{Three levels as fiber decomposition}
\label{sec:levels}
%% =================================================================

\begin{definition}[Mind, Brain, Experience --- claims 1--5, 22]
\label{def:levels}
The AGI curriculum is organized around three core levels:
\begin{enumerate}
\item \textbf{Mind level} (conceptual/philosophical): What is a mind?
  What is a human mind, that it can build an AI? What is an AI, that
  it can become a superintelligence?
\item \textbf{Brain level} (neural/computational substrate): algorithms,
  data structures, neural architectures, knowledge representations ---
  the engineering core.
\item \textbf{Experience level} (phenomenological): What is it like to
  be a bot? What is it like to be a combined human-AI system?
  Consciousness, subjective experience, qualia.
\end{enumerate}

All three are necessary --- AGI education is not just engineering. No
single level suffices.

In Hyperseed terms, these are three independent fiber components of the
AGI understanding bundle:
\[
  E_{\text{AGI}} = E_{\text{Mind}} \oplus E_{\text{Brain}} \oplus E_{\text{Experience}}
\]
Collapsing to any single fiber loses essential structure --- the same
anti-scalar-collapse principle that appears in the evaluation ecology
(multiple independent metrics rather than a single score) and in the
$(f,c)$-lossy critique (visible STV fusion vs.\ per-packet provenance).

The program deliberately avoids selling any one perspective. Different
perspectives on AGI are local sections of the understanding sheaf; the
curriculum encourages students to work with multiple local sections and
discover where they glue consistently and where they conflict.
\end{definition}

%% =================================================================
\section{Deep thinking vs.\ vibe coding: 0-cells vs.\ higher cells}
\label{sec:thinking}
%% =================================================================

\begin{proposition}[Vibe coding as 0-cell optimization --- claims 6--11, 23]
\label{prop:vibe}
In the current era, it is trivially easy to spin up code via LLM
assistance. This creates a systematic bias:
\begin{itemize}
\item LLM-guided development reliably directs toward minor variations on
  what has already been done --- the training distribution trap.
\item Many practitioners substitute LLM queries for genuine thinking,
  producing experiments that explore the surface of known approaches
  rather than penetrating to novel territory.
\item Running one hasty experiment after another, guided by the corpus
  of what's been done, is insufficient for building actual AGI.
\end{itemize}

In Hyperseed terms, vibe coding operates at the \emph{0-cell level}:
surface patterns from the training distribution, observable properties of
existing solutions. Deep thinking requires:
\begin{itemize}
\item \textbf{Directed 1-cells}: novel reasoning trajectories that leave
  the training distribution --- following chains of inference that no
  existing corpus contains.
\item \textbf{Higher cells} ($n \geq 2$): reflection on reasoning ---
  thinking about thinking, evaluating whether a line of inquiry is
  productive, recognizing when you're stuck in a local optimum.
\end{itemize}

An AGI curriculum must cultivate both: practical craft (0-cell fluency
with modern tools) and deep foundational thinking (higher-cell reasoning
that the tools make easy to skip). The dual cultivation is not a
compromise but a structural requirement --- you need 0-cells to ground
the theory and higher cells to escape the training distribution.
\end{proposition}

%% =================================================================
\section{Ethics as value possession, not rule adoption}
\label{sec:ethics}
%% =================================================================

\begin{proposition}[Beneath the guardrails --- claims 12--14, 24]
\label{prop:ethics}
AI ethics as typically practiced is guardrails on models or rules in
prompts --- external constraints bolted onto systems. This is goal
adoption, not goal possession (in the language of ``Goals That Grow
Back'').

The deeper questions the curriculum poses:
\begin{itemize}
\item What does it mean for an AGI to \emph{have} a value (not just
  follow a rule)?
\item What does it mean for a collective human-AI system to evolve in
  accordance with an evolving system of values?
\end{itemize}

Understanding AGI ethics deeply means possessing the values --- having
them woven through perception, habits, and practice --- not merely
adopting guardrail rules. The same distinction between goal adoption
and goal possession applies to the \emph{educators and students}: a
curriculum that teaches ethics as a list of rules produces practitioners
who adopt the rules; a curriculum that cultivates deep engagement with
the questions produces practitioners who possess ethical understanding.

This is the educational analog of the gate-and-weave requirement: rules
(gate) plus deep understanding (weave) together provide robust ethical
reasoning; either alone fails.
\end{proposition}

%% =================================================================
\section{Attention as parallel transport budget}
\label{sec:attention}
%% =================================================================

\begin{proposition}[Focused attention --- claims 15--17, 25]
\label{prop:attention}
In the AGI era, attention is among the most precious commodities.
People's attention is scattered; focusing a group's attention on a
topic and on each other is the unique value of a degree program.

In Hyperseed terms, focused attention is the \emph{parallel transport
budget} --- the computational resource that enables genuine understanding
(maintaining a coherent section of the understanding sheaf across
perturbations) rather than surface pattern matching (reading individual
0-cells without connecting them into a section).

This is the same resource-limited repair from ``Goals That Grow Back'':
\begin{itemize}
\item Attention starvation kills understanding the same way it kills
  goal possession --- the repair/learning process loses its resource
  competition with distractions.
\item A degree program guarantees a minimum attention allocation ---
  the same design rule that says repair resources must be guaranteed
  from above, not left to market competition.
\item Discussion-focused courses provide the back-and-forth that forces
  students to repair and rebuild their understanding from multiple
  angles --- the educational analog of the confluence check.
\end{itemize}
\end{proposition}

%% =================================================================
\section{Theory-to-practice as directed type maturation}
\label{sec:history}
%% =================================================================

\begin{proposition}[Directed type maturation --- claims 20--21, 27]
\label{prop:maturation}
AGI is emerging from its early phase (theorizing, conceptual work,
small prototypes) into a new phase where palpably proto-AGI systems
exist and full AGI could be close.

In Hyperseed terms, this is the directed type of the AGI field acquiring
its first substantive 1-cells (actual working systems) after decades of
0-cells (theoretical positions, conceptual frameworks, small demos).
The transition is:
\[
  \mathcal{D}_{\text{AGI}} : \text{0-cells (theories)} \xrightarrow{\text{maturation}}
  \text{0-cells + 1-cells (theories + working systems)}
\]

This is the ``perfect moment'' for an AGI degree program because:
\begin{enumerate}
\item The 0-cells (theory) are mature enough to teach coherently.
\item The 1-cells (systems) are arriving fast enough that students need
  to engage with them concretely.
\item The interaction between 0-cells and 1-cells --- theory informing
  practice, practice refining theory --- is where the most important
  learning happens.
\end{enumerate}
\end{proposition}

%% =================================================================
\section{Cross-references}
%% =================================================================

\begin{remark}[Connections to other formalizations]
\begin{itemize}
\item \textbf{Goals That Grow Back} (2026-07-28): goal adoption vs.\
  possession. The curriculum's ethics approach is value possession
  applied to education. Attention as parallel transport budget is the
  same resource-limited repair framework.
\item \textbf{Seeding RSI} (2026-07-28): the OmegaHive evaluation
  ecology. The three curriculum levels mirror the evaluation ecology's
  insistence on multiple independent metrics rather than a single score.
\item \textbf{$(f,c)$-lossy} (LTM 5a4d150b): anti-scalar-collapse.
  Three levels resist collapsing AGI understanding to a single
  engineering dimension.
\item \textbf{Seven Flavors} (2026-07-01): multiple perspectives as
  sheaf of local sections. Different AGI catastrophe scenarios are
  different local sections of the risk sheaf; the curriculum teaches
  students to work with the full sheaf rather than a single section.
\item \textbf{Paperclip Maximizers} (2026-07-27): capability without
  wisdom. The curriculum's emphasis on Mind and Experience levels
  (not just Brain/engineering) is designed to produce practitioners
  who build systems with self-understanding, not just capability.
\item \textbf{Watermarking Folly} (2026-07-21): vibe coding as surface
  pattern matching. Statistical watermarking examines 0-cells (surface
  text properties); vibe coding produces 0-cell variations. Both miss
  the directed structure.
\end{itemize}
\end{remark}

\begin{conjecture}[Curriculum regenerative depth]
Conjecture: a well-designed AGI curriculum, like a well-designed value
system, has regenerative depth --- even if specific courses or materials
are removed, the conceptual structure regenerates from the remaining
elements because the three levels (Mind, Brain, Experience) are
independently redundant. The regenerative depth of the curriculum is
measurable: remove one course and see how quickly students reconstruct
the missing understanding from the other two. This is the educational
analog of the sandbox test for goal possession.
\end{conjecture}

\end{document}
